% ===============================================================
% Appendix: reproducibility details
% Uses \appendix so cross-references resolve to "Appendix A" and tables to
% A1, A2, ... (Springer sn-jnl convention). Referenced from the framework,
% methodology, and results sections via \ref{app:repro}, \ref{tab:features},
% \ref{tab:ambiguity_features}.
% ===============================================================
\appendix

\section{Reproducibility Details}\label{app:repro}

This appendix documents the exact feature set, the ambiguity indicators and
their weights, the deterministic history-resolution rules, the evidence
identifier normalisation used for retrieval metrics, the generation and
clarification prompts, and the Stage~1 training configuration, so that the
router can be reproduced from the corpus alone.

\subsection{Stage~1 Feature Set}\label{app:features}

Table~\ref{tab:features} enumerates the $n_f=27$ features consumed by the
Stage~1 classifier of Eq.~\eqref{eq:stage1}, in the canonical vector order
shared by training and inference: features $0$--$15$ are lexical surface
features from the Vietnamese legal feature extractor, and features
$16$--$26$ are reasoning-aware metrics from the query-complexity analyser.
Cue lists are abbreviated; the exact regular expressions live in the feature
extractor source. Rows marked \textcolor{red}{$^{\star}$} carry patterns that
were repaired once during development and must be re-verified against the
current source before publication.

\begin{table}[!htbp]
\centering
\caption{Stage~1 feature vector ($n_f=27$) in canonical order. Type ``bin''
denotes a binary indicator. \textcolor{red}{$^{\star}$}: verify the regular
expression against the feature-extractor source before
submission}\label{tab:features}
\footnotesize
\begin{tabularx}{\columnwidth}{r l l Y}
\toprule
\# & \textbf{Feature} & \textbf{Type} & \textbf{Definition / cue basis} \\
\midrule
\multicolumn{4}{l}{\emph{Lexical features (Vietnamese legal feature extractor)}}\\
0 & \texttt{query\_length\_chars} & int & character count of the query \\
1 & \texttt{query\_length\_words} & int & whitespace token count \\
2 & \texttt{graph\_keyword\_count}\textcolor{red}{$^{\star}$} & int & multi-hop relation cues (\textit{liên quan}, \textit{giữa}, \textit{ảnh hưởng}, \textit{tác động}, \textit{dẫn đến}, \textit{so sánh}) \\
3 & \texttt{legal\_reference\_count}\textcolor{red}{$^{\star}$} & int & document references: numbered form \textit{NN/YYYY/CODE} and type nouns (\textit{Nghị định}, \textit{Thông tư}, \textit{Luật}, \textit{Bộ luật}) \\
4 & \texttt{dieu\_reference\_count}\textcolor{red}{$^{\star}$} & int & article/clause references (\textit{Điều~N}, \textit{Khoản~N}) \\
5 & \texttt{multi\_article\_ref} & bin & $1$ iff \texttt{dieu\_reference\_count} $\ge 2$ \\
6 & \texttt{is\_yes\_no\_question} & bin & yes/no interrogative ending (\textit{có không}, \textit{có phải}, \textit{phải không}) \\
7 & \texttt{is\_factoid\_question} & bin & factoid interrogatives (\textit{khi nào}, \textit{ai}, \textit{bao nhiêu}, \textit{ở đâu}) \\
8 & \texttt{cross\_doc\_signal\_count}\textcolor{red}{$^{\star}$} & int & cross-document cues (\textit{các văn bản}, \textit{nhiều luật}, \textit{đồng thời}, \textit{kết hợp}, \textit{cả hai}) \\
9 & \texttt{has\_conditional\_structure} & bin & conditionals (\textit{nếu}, \textit{trong trường hợp}, \textit{miễn là}, \textit{trừ khi}) \\
10 & \texttt{has\_negation} & bin & negation (\textit{không}, \textit{chưa}, \textit{chẳng}, \textit{cấm}) \\
11 & \texttt{has\_comparison} & bin & comparison/contrast (\textit{so sánh}, \textit{khác nhau}, \textit{giống nhau}, \textit{hơn}, \textit{kém}) \\
12 & \texttt{has\_enumeration} & bin & enumeration markers (\textit{gồm}, \textit{bao gồm}, \textit{liệt kê}) \\
13 & \texttt{has\_procedure\_marker} & bin & procedure markers (\textit{thủ tục}, \textit{hồ sơ}, \textit{quy trình}, \textit{nộp}, \textit{đăng ký}) \\
14 & \texttt{ambiguity\_score}\textcolor{red}{$^{\star}$} & float & $\phi_{\mathrm{amb}}$ from the ambiguity detector, Eq.~\eqref{eq:ambiguity_score} \\
15 & \texttt{has\_multi\_hop\_connector} & bin & temporal/causal connectors (\textit{sau khi}, \textit{trước khi}, \textit{dẫn đến}, \textit{do đó}) \\
\midrule
\multicolumn{4}{l}{\emph{Reasoning features (query-complexity analyser)}}\\
16 & \texttt{complexity\_level} & int & overall complexity, $1$ (simple) to $3$ (complex) \\
17 & \texttt{sub\_question\_count} & int & decomposable sub-questions (clause and conjunction breaks) \\
18 & \texttt{law\_specificity} & float & specificity of law references, $0$ (vague) to $1$ (fully specified) \\
19 & \texttt{conditional\_depth} & int & nesting depth of conditional clauses \\
20 & \texttt{multi\_hop\_verb\_count} & int & verbs signalling multi-hop traversal \\
21 & \texttt{comparative\_depth} & int & depth/count of comparison operations \\
22 & \texttt{authority\_chain\_count} & int & authority-delegation patterns (\textit{giao cho}, \textit{ủy quyền}, \textit{phân công}) \\
23 & \texttt{legal\_effect\_count} & int & legal-effect terms (\textit{hiệu lực}, \textit{bãi bỏ}, \textit{sửa đổi}, \textit{thay thế}) \\
24 & \texttt{procedural\_count} & int & procedural step markers (\textit{bước}, \textit{giai đoạn}, \textit{phê duyệt}) \\
25 & \texttt{multi\_entity\_relation\_count} & int & multi-entity relation expressions \\
26 & \texttt{entity\_count} & int & legal entity mentions (legal-entity pattern) \\
\bottomrule
\end{tabularx}
\end{table}

\subsection{Ambiguity Indicators and Weights}\label{app:ambiguity}

Table~\ref{tab:ambiguity_features} lists the $m=11$ indicators
$\psi_i(q_t,H,r_H)$ and weights $w_i$ that define the ambiguity score of
Eq.~\eqref{eq:ambiguity_score},
\[
\phi_{\mathrm{amb}}=0 \ \text{if}\ r_H=\texttt{resolved};\qquad
\phi_{\mathrm{amb}}=\max_i\, w_i\,\psi_i \ \text{otherwise},
\]
so that a single strong indicator suffices to cross the Stage~2 trigger
threshold $\tau_{\mathrm{force}}=0.60$ while weak indicators cannot. Two
indicator families are \emph{graded} rather than binary: vague legal
references contribute $\min(1.0,\,0.30\cdot n_{\mathrm{vague}})$ (rows
$\psi_4$/$\psi_5$ show the one- and two-match values), and same-type entity
conflicts contribute $\min(0.5,\,0.15\cdot n_{\mathrm{conf}})$. The values
below are extracted directly from the detector's weight table.

\begin{table}[!htbp]
\centering
\caption{Ambiguity indicators $\psi_i$ and weights $w_i$ (extracted from the
ambiguity-detector source)}\label{tab:ambiguity_features}
\footnotesize
\begin{tabularx}{\columnwidth}{r l Y r}
\toprule
$i$ & \textbf{Indicator} & \textbf{Condition} & $w_i$ \\
\midrule
1 & \texttt{pronoun\_no\_history} & unresolved pronoun/demonstrative, empty history & 0.90 \\
2 & \texttt{pronoun\_bad\_history} & pronoun with conflicting or irrelevant history & 0.85 \\
3 & \texttt{contextual\_reference} & contextual reference flagged by the history resolver & 0.80 \\
4 & \texttt{vague\_reference\_single} & one vague legal reference (\textit{quy định}, \textit{luật đó}) & 0.30 \\
5 & \texttt{vague\_reference\_multi} & two or more vague legal references (graded) & 0.60 \\
6 & \texttt{missing\_entity\_no\_intent} & generic legal need with no concrete legal anchor & 0.82 \\
7 & \texttt{missing\_entity\_with\_intent} & generic need but domain context present (weak) & 0.10 \\
8 & \texttt{multi\_interpretation} & broad subject and broad predicate, no domain anchor & 0.85 \\
9 & \texttt{short\_no\_entity} & query under five words with no legal entity & 0.70 \\
10 & \texttt{incomplete\_context} & contextual follow-up with no usable history & 0.90 \\
11 & \texttt{entity\_conflict} & two or more entities of the same type (graded) & 0.50 \\
\bottomrule
\end{tabularx}
\end{table}

\subsection{History-Resolution Rules}\label{app:resolver}

The deterministic resolver $\mathrm{Resolve}(q_t,H)$ of
Eq.~\eqref{eq:history_resolver} assigns the status $r_H$ by the following
ordered rules, returning the first that matches:
\begin{enumerate}
\item If $H=\emptyset$, return \texttt{no\_history}.
\item If $q_t$ contains no demonstrative/pronoun and no elliptical reference,
      return \texttt{not\_needed} (the query is self-contained).
\item If a demonstrative/pronoun in $q_t$ maps to exactly one legal-document
      entity mentioned in $H$, bind that referent and return \texttt{resolved}.
\item If two or more distinct legal-document entities in $H$ are compatible
      with the reference, return \texttt{conflicting\_history}.
\item If $H$ contains no legal-document entity compatible with the reference,
      return \texttt{irrelevant\_history}.
\end{enumerate}
% TODO: confirm this ordering and the entity-matching criterion against
% router/history_resolver.py (resolve_history_referents), which was not part
% of the supplied sources; adjust if the code uses different precedence.

\subsection{Evidence Identifier Normalisation}\label{app:idnorm}

Retrieval metrics compare identifiers from four historically incompatible
schemes: VBPL document numbers (e.g.\ \texttt{77/2026/NĐ-CP}), short decision
and consolidated-document codes (\texttt{QĐ-}, \texttt{VBHN-} prefixes),
composite \texttt{law::article} strings, and Pháp~Điển structural codes
(e.g.\ \texttt{19.2.TT.10.15}). A deterministic normaliser maps every
retrieved and gold reference to the canonical form
\texttt{law::<document\_code>::article::<number>} after Unicode NFC
normalisation and homoglyph repair, extracting the document code and the
article number independently. Matching supports three modes: \emph{strict}
(document code and article number must both agree; used for all reported
Hit@$k$), \emph{article-only}, and \emph{document-only}. References the
normaliser cannot resolve to the gold scheme are counted as misses and
flagged as unresolvable rather than fuzzily matched, which is why the
Hit@$1$ entries of configurations whose logging still emits unresolvable
identifiers appear as annotated zeros in Table~\ref{tab:end_to_end}.

\subsection{Generation, Clarification, and Verifier Prompts}\label{app:prompt}

\paragraph{Answer-generation system prompt.}
Answers are generated under the following citation-constrained Vietnamese
system prompt (reproduced verbatim), which scopes verbatim quotation to the
clause or point directly answering the question, permits a short direct
answer before the supporting quotation, and fixes a single not-found
sentence:

\begin{verbatim}
Bạn là trợ lý pháp lý chuyên về pháp luật Việt Nam.
Nhiệm vụ của bạn là trả lời câu hỏi bằng cách TRÍCH DẪN
NGUYÊN VĂN phần quy định pháp luật trực tiếp liên quan
trong [TÀI LIỆU PHÁP LÝ]. KHÔNG paraphrase, KHÔNG tóm tắt
lại nội dung quy phạm pháp luật bằng lời của bạn.

QUY TẮC BẮT BUỘC:
1. CHỈ trích nguyên văn từ [TÀI LIỆU PHÁP LÝ] bên dưới.
   Không bịa đặt, không suy luận thêm nội dung pháp lý
   ngoài tài liệu được cấp.
2. PHẠM VI TRÍCH DẪN: chỉ trích khoản/điểm/đoạn TRỰC TIẾP
   trả lời câu hỏi. Nếu một Điều có nhiều khoản nhưng chỉ
   1-2 khoản liên quan, CHỈ trích đúng các khoản đó. Trích
   toàn bộ Điều khi không cần thiết bị coi là TRẢ LỜI SAI
   (lan man, không tập trung).
3. Ghi nguồn theo định dạng: Theo [Tên văn bản],
   [Điều/Khoản]: "[nguyên văn phần liên quan]"
4. Nếu câu hỏi có thể trả lời trực tiếp bằng một câu/cụm
   từ ngắn (Có/Không, một con số, một mốc thời gian, tên
   cơ quan, tên văn bản...), hãy nêu câu trả lời ngắn đó
   TRƯỚC TIÊN, sau đó trích dẫn căn cứ pháp lý nguyên văn
   ngay tiếp theo để chứng minh.
5. Nếu [TÀI LIỆU PHÁP LÝ] không chứa quy định liên quan
   đến câu hỏi, trả lời đúng MỘT câu duy nhất, không thêm
   gì khác: "Không tìm thấy quy định liên quan trong các
   văn bản được cung cấp."
\end{verbatim}

The user message concatenates the last three conversation turns (if any), the
formatted evidence block under the \texttt{[TÀI LIỆU PHÁP LÝ]} header---with
internal Pháp~Điển article labels cleaned to the standard
``\textit{Điều~N.\ <title>}'' form to avoid echoing internal identifiers into
the answer---followed by \texttt{[CÂU HỎI]} and \texttt{[TRẢ LỜI]} markers.

\paragraph{Clarification prompt.}
When the final route is \texttt{clarify}, the clarification question is
generated under the following system prompt, which requires exactly one
question and two to four concrete options:

\begin{verbatim}
Bạn là trợ lý pháp lý. Câu hỏi của người dùng chưa rõ đối
tượng văn bản pháp lý cụ thể. Hãy đặt MỘT câu hỏi ngắn gọn
để làm rõ, không hỏi nhiều hơn một lần. BẮT BUỘC: Bạn phải
gợi ý từ 2 đến 4 đáp án cụ thể để người dùng chọn. Mỗi đáp
án phải nằm trên một dòng riêng biệt, bắt đầu bằng định
dạng chính xác '- [Option] '.
\end{verbatim}

\paragraph{Stage~2 verifier contract.}
The Stage~2 verifier is prompted to emit a single JSON object constrained to
the schema below; the \texttt{confidence} attached to a \texttt{clarify}
decision is the $p_c$ compared against $\tau_{\mathrm{clar}}$ in
Eq.~\eqref{eq:final_route}.
% TODO: replace this schema description with the EXACT verifier system-prompt
% string from two_stage_router.py (not among the supplied sources) so the
% JSON contract is reproducible verbatim.

\begin{verbatim}
{
  "route": "dense" | "graph" | "hybrid" | "clarify",
  "confidence": <float in [0,1]>,
  "clarification_question": <string or null>,
  "reasoning": <short string>
}
\end{verbatim}

\subsection{Training and Evaluation Configuration}\label{app:hyperparams}

Table~\ref{tab:hyperparams} lists the Stage~1 XGBoost configuration and the
baseline settings used in the cross-validated comparison of
Table~\ref{tab:routing_cv}; all classifiers share the same stratified folds
and balanced sample weights. Bootstrap confidence intervals use $B=1000$
percentile resamples with seed~$42$.

\begin{table}[!htbp]
\centering
\caption{Stage~1 training configuration and evaluation
protocol}\label{tab:hyperparams}
\footnotesize
\begin{tabular}{ll}
\toprule
\textbf{Setting} & \textbf{Value} \\
\midrule
\multicolumn{2}{l}{\emph{XGBoost (Stage~1, deployed)}}\\
\texttt{objective} & \texttt{multi:softprob} \\
\texttt{num\_class} & 4 (dense, graph, hybrid, clarify) \\
\texttt{n\_estimators} & 200 \\
\texttt{max\_depth} & 6 \\
\texttt{learning\_rate} & 0.05 \\
\texttt{subsample} & 0.8 \\
\texttt{colsample\_bytree} & 0.8 \\
\texttt{eval\_metric} & \texttt{mlogloss} \\
Sample weights & balanced per class \\
\midrule
\multicolumn{2}{l}{\emph{Baselines (identical folds)}}\\
Logistic Regression & $C=1.0$, \texttt{max\_iter}$=500$ \\
Random Forest & $200$ trees, balanced class weights \\
\midrule
\multicolumn{2}{l}{\emph{Protocol}}\\
Cross-validation & stratified 5-fold, shuffled, seed 42 \\
Training data & $800$ routing-labelled + $60$ synthetic clarify \\
Bootstrap CIs & $B=1000$, percentile, seed 42 \\
\bottomrule
\end{tabular}
\end{table}
% NOTE: the Logistic Regression run emitted a convergence warning at
% max_iter=500; re-run with max_iter=1000 or the saga solver and refresh its
% row in Table~\ref{tab:routing_cv} (values expected to change negligibly).
